\documentclass{article}
\usepackage{/globals/anki}
\ankiprefix{wi-software}
\ankitopic{WI Software}
\ankishow
\begin{document}
\section{Software} 
Software kann in zwei Arten unterschieden werden; Standardsoftware und Individualsoftware. \ankinote{arten}{Arten}{Standard, Individual}

\subsection{Standardsoftware}
Standardsoftare wird von anderen Firmen erstellt und in Massen verkauft. \ankinote{standard}{Standard}{In Massen verkauft}
\begin{description}
 \item[Basissoftware] Alltägliche Software, wie Webbrowser oder Antivirensoftware.
 \item[Standardbürosoftware] Software des Büroalltags, wie Officeprogramme oder Kommunikationssoftware.
 \item[Funktionsorientierte Software] Software die eine bestimmte betriebswirtschaftliche Funktion hat.
 \item[Prozessorientierte Software] Funktionsübergreifende Software um Prozesse abzubilden.
\end{description}
\ankinote{standard-arten}{Standard Arten}{Basis, Standardbüro, Funktionsorientiert, Prozessorientiert}
\ankinote{basis}{Basis}{Zuhause Alltag}
\ankinote{standardbüro}{Standardbüro}{Büro Alltag}
\ankinote{funktionsorientiert}{Funktionsorientiert}{Spezifische betriebswirtschaftliche Funktion}
\ankinote{prozessorientiert}{Prozessorientiert}{Übergreifende Funktion, Prozess}
\subsection{Invidivualsoftware}
Invidivualsoftware wurde spezifisch für die bestimmten eigenen Anforderungen entwickelt. Somit ist es um einiges teurer, differenziert einen aber auch vom Wettbewerb.
\ankinote{individual}{Individual}{spezifisch, teuer, einzigartig}

Der Trend geht eher von der Indivualsoftware weg.
\ankinote{individual-trend}{Trend}{Zu Standard}
\end{document}